% !TeX program = xelatex
% Standalone report package. Compile from this directory; figures and tables are local.
\documentclass[UTF8,a4paper,zihao=-4,fontset=none]{ctexart}
\usepackage[margin=24mm,headheight=18pt]{geometry}
\usepackage{fontspec}
\IfFontExistsTF{FangSong}{\setCJKmainfont{FangSong}[BoldFont=SimHei]\setCJKsansfont{SimHei}}{\setCJKmainfont{FandolFang-Regular.otf}[BoldFont=FandolHei-Bold.otf]\setCJKsansfont{FandolHei-Regular.otf}}
\setmainfont{texgyretermes-regular.otf}[BoldFont=texgyretermes-bold.otf,ItalicFont=texgyretermes-italic.otf,BoldItalicFont=texgyretermes-bolditalic.otf]
\setsansfont{texgyreheros-regular.otf}[BoldFont=texgyreheros-bold.otf]
\setmonofont{lmmono10-regular.otf}[Scale=MatchLowercase]
\usepackage{amsmath,amssymb,graphicx,booktabs,tabularx,longtable,array,multirow}
\usepackage[table]{xcolor}
\usepackage{setspace,caption,float,placeins,enumitem,fancyhdr,titlesec,needspace,etoolbox}
\usepackage[unicode,hidelinks]{hyperref}
\usepackage{xurl,listings}
\definecolor{Navy}{HTML}{17334D}
\definecolor{Teal}{HTML}{167D8D}
\definecolor{Pale}{HTML}{EEF4F7}
\definecolor{Gray}{HTML}{546572}
\definecolor{Gold}{HTML}{A56E28}
\definecolor{LightGray}{HTML}{F4F6F8}
\newcolumntype{Y}{>{\raggedright\arraybackslash}X}
\newcolumntype{L}[1]{>{\raggedright\arraybackslash}p{#1}}
\newcommand{\code}[1]{{\small\nolinkurl{#1}}}
\newcommand{\tabnote}[1]{\par\vspace{3pt}{\footnotesize\color{Gray}注：#1\par}}
\newcommand{\fig}[4]{\begin{figure}[!htbp]\centering\includegraphics[width=#3\linewidth]{figures/#1.pdf}\caption{#2}\label{#4}\end{figure}}
\newcommand{\status}[1]{\leavevmode\textcolor{Teal}{#1}}
\newcommand{\pending}[1]{\leavevmode\textcolor{Gold}{#1}}
\setstretch{1.5}
\makeatletter
\renewcommand\normalsize{\@setfontsize\normalsize{14bp}{16.8bp}}
\makeatother
\setlength{\parindent}{2em}
\setlength{\parskip}{2pt}
\setlength{\emergencystretch}{2em}
\setlength{\tabcolsep}{5pt}
\renewcommand{\arraystretch}{1.22}
\captionsetup{font=small,labelfont=bf,labelsep=quad,skip=8pt}
\captionsetup[table]{position=top}
\setlist{nosep,leftmargin=2em}
\ctexset{section={format=\sffamily\bfseries\zihao{3}\color{Navy},name={,},number=\chinese{section},aftername=、,beforeskip=22pt,afterskip=10pt},subsection={format=\sffamily\bfseries\zihao{4}\color{Navy},beforeskip=16pt,afterskip=6pt},subsubsection={format=\bfseries\zihao{4},beforeskip=12pt,afterskip=4pt},contentsname={目\quad 录}}
\pagestyle{fancy}
\fancyhf{}
\fancyhead[L]{\footnotesize\sffamily\color{Gray}智能科研助理系统\enspace /\enspace 模块五评测报告}
\fancyhead[R]{\footnotesize\sffamily\color{Gray}成员 B}
\fancyfoot[C]{\small\thepage}
\renewcommand{\headrulewidth}{0.3pt}
\lstset{basicstyle=\ttfamily\footnotesize,breaklines=true,columns=fullflexible,keepspaces=true,backgroundcolor=\color{LightGray},frame=single,rulecolor=\color{Pale},xleftmargin=0pt,aboveskip=9pt,belowskip=9pt}
\pretocmd{\section}{\FloatBarrier\Needspace{7\baselineskip}}{}{}
\pretocmd{\subsection}{\Needspace{7\baselineskip}}{}{}
\AtBeginEnvironment{table}{\setstretch{1.1}}
\AtBeginEnvironment{longtable}{\setstretch{1.1}}
\hypersetup{pdftitle={智能科研助理系统：模块五评测与交付报告},pdfauthor={成员B（姓名待填）},pdfsubject={南京农业大学课程实践；55题Agent基线与495条检索实验}}

\begin{document}
\hypersetup{pageanchor=false}
\begin{titlepage}
\thispagestyle{empty}
\noindent{\sffamily\color{Navy}\zihao{2}\bfseries 南京农业大学}\par
\noindent{\color{Gray}\zihao{4}智慧农业学院（人工智能学院）\quad 课程实践报告}\par
\vspace{22mm}
\noindent{\color{Teal}\rule{30mm}{2pt}}\par
\vspace{7mm}
\noindent{\sffamily\bfseries\color{Navy}\zihao{1}智能科研助理系统}\par
\vspace{5mm}
\noindent{\sffamily\bfseries\color{Navy}\zihao{2}模块五：系统评测与项目交付}\par
\vspace{8mm}
\noindent{\color{Gray}\large 55 题 Agent 基线\quad /\quad 495 条检索实验\quad /\quad 证据驱动分析}\par
\vspace{19mm}
\noindent\begin{tabularx}{\linewidth}{@{}L{30mm}Y@{}}
\toprule
课程 & 生产实习（1 学分）\\
学年学期 & 2026--2027 学年第 1 学期\\
负责成员 & B：\underline{\hspace{46mm}}（姓名待填）\\
学号 / 班级 & \underline{\hspace{72mm}}\\
指导教师 & \underline{\hspace{72mm}}\\
报告日期 & 2026 年 10 月 14 日\\
\bottomrule
\end{tabularx}
\vfill
\noindent\begin{minipage}{\linewidth}\small\color{Gray}
报告范围：整理成员 B 已完成的语料与评测集准备、实验组织、AI 辅助评分、检索参数分析与错误归因。质量评分尚待独立人工复核；生成/Agent 修复后的质量复测及全项目真实界面截图列为待验收项。姓名、学号等空白项由本人填写。
\end{minipage}
\end{titlepage}

\pagenumbering{Roman}
\hypersetup{pageanchor=true}
\section*{摘\quad 要}
\addcontentsline{toc}{section}{摘要}
本报告依据课程实践方案与项目交付模板，记录成员 B 在模块五的阶段性工作。在去重后的 15 篇、296 个物理页的本地文献上，建立包含事实、对比、总结和推理四类的 55 题评测集，并完成 55 题 Agent 基线及 3 档递归字符分块长度、3 种检索方式组成的 495 条独立检索实验。评测保留失败样本、逐题时延、实际模型调用 Token、工具轨迹和运行来源；使用固定参考答案及原文证据开展 AI 辅助量表评分。

Agent 基线有 53 题执行成功，自动工具集合匹配为 46/55；这两项均不代表答案正确率。正确性、完整性和引用准确性均分分别为 0.9818、0.8182、0.9636（满分 2）。在本次检索实验中，512 字符分块配合混合检索与重排获得 Hit@5=83.64\%、MRR@5=0.6245、Recall@5=72.73\%，平均查询耗时 1.7860 秒，适合作为后续固定基线。错误分析显示，工具选择、论文范围绑定以及证据到回答的忠实映射仍需改进。本报告提供可追溯的数据、矢量图表和复现说明，不将参数筛选或运行恢复误写为已完成的答案质量优化闭环。

\noindent\textbf{关键词：}检索增强生成；Agent；页级证据评测；参数对比；AI 辅助评分；Bad Case
\clearpage
{\ctexset{section/format=\sffamily\bfseries\zihao{2}\color{Navy}}\tableofcontents}
\clearpage
\pagenumbering{arabic}

\section{项目概述}
\subsection{项目背景}
科研文献数量持续增加，人工查找实验配置、比较方法和核验结论耗时较长。本项目构建基于本地论文知识库的智能科研助理，通过检索增强生成与工具调用，支持证据查询、论文比较及归纳推理。成员 B 负责以固定问题集检验检索、回答、引用与运行效率，为系统改进提供可复查依据。

\subsection{目标与报告边界}
模块五的核心目标是形成“数据准备—执行测量—答案审阅—错误分析—改进复测”的可追溯过程。本阶段完成前三项及错误分析，并形成检索参数选择依据与优化交接清单；修改后答案质量复测尚未完成。报告按交付模板的七章结构组织，其他模块只说明评测依赖与可观察状态，完整验收由对应负责人补齐。

评测资料与报告整理使用了 AI 辅助，55 题质量分数来自 AI 对固定参考答案、实际回答、可见工具轨迹及原文页的审阅。它们不是独立人工实验，也不是教师评分。课程要求的人工评估仍需本人或独立评审者复核并签署记录。

\subsection{实际技术栈}
被测系统采用 Python、pdfplumber、Chroma、BM25、RRF 融合、BGE 嵌入与重排、Ollama 本地语言模型，以及 Streamlit 前端。评测层独立记录 JSONL、CSV 和 JSON，再由 Python/Matplotlib 生成统计图，以 LaTeX 组织报告。实验模型与主要配置见附录表\ref{tab:env}，不据此推断未实测的并发能力或部署能力。

\section{团队分工说明}
\begin{table}[H]
\caption{团队约定分工与本报告归属}
\centering\small
\begin{tabularx}{\linewidth}{@{}L{13mm}L{18mm}Y L{25mm}r@{}}
\toprule 成员 & 姓名 & 主要职责 & 对应模块 & 约定比例\\\midrule
A & 待本人填写 & 前端交互、页面呈现和界面材料 & 模块四前端 & 15\%\\
B & 待本人填写 & 评测集、实验运行、评分审阅、错误分析、交付文档 & 模块五 & 15\%\\
C & 待本人填写 & 文档处理、分块、检索与 RAG 生成后端 & 模块一、二 & 35\%\\
D & 待本人填写 & Agent 工具、路由与系统集成 & 模块三、四集成 & 35\%\\\bottomrule
\end{tabularx}
\tabnote{比例来自团队分工方案，合计 100\%，不等同于自动统计的实际代码贡献。本报告不代替其他成员的工作验收。}
\end{table}

成员 B 的实际工作包括：在 AI 辅助下整理与核验 15 篇语料和 55 题参考证据；在本地环境组织模型预检、正式运行与中断恢复；完成 9 组检索参数的结果对照；开展有审计记录的 AI 辅助评分；整理四类 Bad Case 及修复建议；冻结数据并制作本报告图表。检索与 Agent 业务代码分别作为 C、D 提供的被测依赖，不能因 B 运行了评测便将其实现归为 B 的独立开发成果。

\clearpage
\section{各模块交付物清单}
\subsection{模块一：文档处理与检索}
\begin{table}[H]\caption{模块一的评测依赖与验收边界}\small\centering
\begin{tabularx}{\linewidth}{@{}L{35mm}L{25mm}Y@{}}\toprule 项目 & 观察状态 & 依据或待补事项\\\midrule
PDF 加载与索引 & \status{已调用验证} & 15 篇共 296 页；512 档形成 2082 块。\\
Word / TXT 加载 & \pending{待实现或核验} & 当前加载器为 PDF，不宣称已支持所有格式。\\
分块与质量对比 & \status{部分完成} & 已测 recursive 的 256/512/1024 长度；三类分块策略质量比较待补。\\
向量、混合、重排 & \status{已完成本轮} & 三档长度各 55 题 × 3 种检索，共 495 条。\\
嵌入与向量库对比 & \pending{待完成} & 本轮固定 BGE 与 Chroma，未测 m3e/FAISS 对照。\\
增量索引 & \pending{待核验/改进} & 当前建库全量重建；复用已有索引不是增量更新。\\\bottomrule\end{tabularx}
\end{table}
\subsection{模块二：RAG 生成}
\begin{table}[H]\caption{模块二的评测依赖与验收边界}\small\centering
\begin{tabularx}{\linewidth}{@{}L{35mm}L{25mm}Y@{}}\toprule 项目 & 观察状态 & 依据或待补事项\\\midrule
上下文生成与引用 & \status{可调用，质量待改进} & Agent 调用了 RAG 工具；有引用错配和信息遗漏。\\
解码参数实验 & \pending{待完成} & 未完成 temperature/top-p/top-k 的独立质量对照。\\
真实流式与语义缓存 & \pending{未验证} & 前端逐字展示不能作为模型 Token 流式的证据。\\
日志与异常恢复 & \status{部分验证} & 有真实调用、耗时及失败记录；完整降级验收由 C/D 补齐。\\\bottomrule\end{tabularx}
\end{table}
\subsection{模块三：Agent}
\begin{table}[H]\caption{模块三的评测依赖与验收边界}\small\centering
\begin{tabularx}{\linewidth}{@{}L{35mm}L{25mm}Y@{}}\toprule 项目 & 观察状态 & 依据或待补事项\\\midrule
工具调用与 ReAct & \status{已评测基线} & 55 题中 27 题走 ReAct，28 题直接工具路由。\\
工具数量与参数 & \status{部分完成} & 当前 7 类工具；课程 8 类目标及参数契约需负责人核验。\\
错误恢复与论文绑定 & \pending{待修复复测} & F07/F12 抛出类型异常；存在论文范围误匹配。\\
记忆与健康检查 & \pending{本轮未验收} & 本轮关闭跨题记忆；不能推断多轮记忆质量。\\\bottomrule\end{tabularx}
\end{table}
\subsection{模块四：集成与界面}
\begin{table}[H]\caption{模块四的评测依赖与验收边界}\small\centering
\begin{tabularx}{\linewidth}{@{}L{35mm}L{25mm}Y@{}}\toprule 项目 & 观察状态 & 依据或待补事项\\\midrule
前端代码与服务接口 & \status{已交付代码} & 本报告不据文件存在推断全部页面已端到端验收。\\
历史记录与核心交互 & \pending{待界面验收} & 需 A/D 在真实运行界面核验。\\
3--5 张真实界面截图 & \pending{待补齐} & 模板 4.4.2 的全项目要求；本报告数据图不替代界面截图。\\\bottomrule\end{tabularx}
\end{table}
\subsection{模块五：系统评测与交付}
\begin{table}[H]\caption{成员 B 的交付与尚待闭环事项}\small\centering
\begin{tabularx}{\linewidth}{@{}L{37mm}L{23mm}Y@{}}\toprule 交付项 & 状态 & 可复查证据\\\midrule
10--20 篇、50+ QA & \status{已完成} & 15 篇、55 题、四类问题与原文证据。\\
检索质量与参数图表 & \status{已完成本轮} & 9 组配置、495 条结果与逐题指标。\\
Agent 与效率基线 & \status{已完成本轮} & 55 题，保留 2 个执行失败及全部计量字段。\\
质量评分 & \status{AI 辅助完成} & 55 题逐项评分、证据及审计；独立人工复核待补。\\
四类 Bad Case & \status{已完成分析} & 检索/生成/引用/Agent 多标签与代表案例。\\
质量优化前后复测 & \pending{待完成} & 已形成 C/D 交接建议，尚无修复后 55 题结果。\\
技术与使用说明 & \status{本报告已整理} & 设计、复现命令、结果说明与源文件。\\\bottomrule\end{tabularx}
\end{table}
\FloatBarrier

\section{技术实现记录}
\subsection{模块一：文档与索引接口}
本轮从项目论文目录筛选文献，剔除同文重复版本与主题外指南，冻结 PDF 文件名和 SHA-256。物理页码从 1 开始，避免印刷页码与 PDF 顺序页混淆。索引流程采用文本提取、递归字符分块、BGE 嵌入、Chroma 存储，再组装 BM25 和融合检索。512 档复用了经过文本、页码、块 ID 与嵌入配置验证的 2082 个已有块；其建库时间不与新建索引直接比较。

\subsection{模块二：生成与计量接口}
回答链路使用 Ollama 的 \code{qwen2.5:7b}。调用级 Token 计数聚合实际返回的 prompt/completion 字段；单题若存在未完整计量的调用，不将缺失值伪装为零。本轮 55 题 Token 覆盖率为 100\%。预热、索引组装与模型加载耗时不并入逐题问答时延。各工具保留原有解码配置，本报告不将预检温度或某一路生成温度称为全链路统一参数。

\subsection{模块三：Agent 执行与轨迹}
系统根据规则直接调用工具，或进入 ReAct 循环。评测记录实际工具集合、参数、返回状态、执行模式和步数。\code{success=True} 仅表示后端声明执行成功；S10 等题可能只返回诊断性文本，因此必须再检查回答内容。自动工具集合匹配与 AI 对参数及任务意图的路由判断分开统计，避免把两个定义不同的指标混为一个准确率。

\subsection{模块四：集成与截图材料}
本阶段通过脚本调用实际后端完成基线。图\ref{fig:architecture}为代码依赖及评测数据流示意，不是界面截图。模板要求的真实截图应由 A/D 在系统运行时补充，建议覆盖文献加载、带引用问答、论文比较、工具轨迹及历史记录，并标注核心区域。本报告没有将模拟网页或数据统计图充作该项证据。

\fig{08_architecture}{被测系统与成员 B 评测流水线。架构示意图；实线表示调用或数据传递，评测层负责冻结输入、计量与审阅。}{1}{fig:architecture}
\FloatBarrier

\subsection{模块五：评测集与实验设计}
\subsubsection{数据集构建与覆盖}
语料为视觉表征、语义分割及视觉应用研究，另纳入作为 Transformer 理论基础的 \emph{Attention Is All You Need}，不将 15 篇都称为视觉应用论文。构造事实 20 题、对比 15 题、总结 10 题、推理 10 题。每题保存 ID、问题、参考答案、证据 PDF/物理页/摘录、可接受工具集合；共核验 142 条证据记录。论文题名与页数见附录表\ref{tab:corpus}。

为保持任务独立性，本轮关闭跨题会话记忆。Agent 各题只执行一次正式基线；三档分块各完成向量、混合、混合加重排三种检索。所有检索配置使用相同 55 题、相同语料与相同后端指纹。图\ref{fig:dataset}中的论文覆盖按题目引用的金标证据去重统计；一题可覆盖多篇论文，因此论文列的计数之和可以大于 55。
\fig{01_dataset_profile}{评测集结构与论文证据覆盖。四类问题共 55 题；P01--P15 的题名对应附录清单。覆盖数量表示关联题数，不是独立问题总数。}{1}{fig:dataset}

\subsubsection{检索指标定义}
设第 $i$ 题的去重金标证据页集合为 $G_i$，前 5 个检索块对应页的集合为 $P_i^{(5)}$，$r_i$ 为前 5 名中首次命中金标页的块排名。证据单位为“精确 PDF 文件名 + 物理页”。
\begin{align}
\mathrm{Hit@5}&=\frac1N\sum_{i=1}^{N}\mathbf{1}[G_i\cap P_i^{(5)}\neq\varnothing],\\
\mathrm{MRR@5}&=\frac1N\sum_{i=1}^{N}\begin{cases}1/r_i,&r_i\leq5,\\0,&\text{前五名无命中},\end{cases}\\
\mathrm{Recall@5}&=\frac1N\sum_{i=1}^{N}\frac{|G_i\cap P_i^{(5)}|}{|G_i|}.
\end{align}
重复同页块不增加 Recall 的已覆盖证据页数，MRR 仍按块的原始排名计算。页级命中不证明整页已进入模型上下文，也不证明答案结论正确。Agent 汇总引用可能来自多次工具调用，没有统一 Top-5 排名，因而其检索三指标为 N/A；下文检索成绩全部来自独立 retrieval 模式。

\subsubsection{固定条件与测量边界}
本轮均使用 recursive 字符分块，长度参数为 256/512/1024，overlap 参数为 50，候选召回 20、最终 Top-5。当前实现不保证每一对输出块都有恰好 50 字符重叠，因此不把配置值解释为实际重叠量保证。嵌入及重排在 CPU 执行，语言模型由 Ollama 管理设备；本报告没有连续 GPU 利用率测量。

55 题 Agent 基线属于中断后续跑：先前已落盘的 23 题保留原始回答、计时与 Token，后续运行执行余下 32 题，各题记录来源。早期异常运行、预检和 3 题冒烟结果不计入正式 55 题。不能用最后一段进程起止时间代表整轮耗时。三档检索都是单次顺序运行，没有重复采样、误差棒或显著性检验。

\subsection{检索参数实验结果}
\begin{table}[H]\caption{九组检索配置的总体表现（每组 $n=55$）}\label{tab:retrieval}\centering\small
% Generated by scripts/make_tables.py from data/report_data.json.
\begingroup
\small
\setlength{\tabcolsep}{5pt}
\renewcommand{\arraystretch}{1.16}
\begin{tabularx}{\linewidth}{@{}r r >{\raggedright\arraybackslash}X r r r r r@{}}
\toprule
\shortstack{块长\\（字符）} & \shortstack{文本块\\（个）} & 检索方法 & \shortstack{命中\\（题/55）} & \shortstack{Hit@5\\（\%）} & MRR@5 & \shortstack{Recall@5\\（\%）} & \shortstack{平均查询\\（s/题）} \\
\midrule
256 & 4\,335 & 向量 & 38/55 & 69.09 & 0.550 & 57.58 & 0.163 \\
 &  & 混合（RRF） & 40/55 & 72.73 & 0.575 & 59.24 & 0.192 \\
 &  & 混合＋重排 & 42/55 & 76.36 & 0.591 & 64.24 & 1.312 \\
\addlinespace[5pt]
512 & 2\,082 & 向量 & 39/55 & 70.91 & 0.576 & 59.85 & \textbf{0.142} \\
 &  & 混合（RRF） & 39/55 & 70.91 & 0.587 & 58.03 & 0.148 \\
 &  & 混合＋重排 & \textbf{46/55} & \textbf{83.64} & \textbf{0.625} & \textbf{72.73} & 1.786 \\
\addlinespace[5pt]
1024 & 1\,079 & 向量 & 40/55 & 72.73 & 0.550 & 57.12 & 0.159 \\
 &  & 混合（RRF） & 42/55 & 76.36 & 0.610 & 61.97 & 0.164 \\
 &  & 混合＋重排 & \textbf{46/55} & \textbf{83.64} & 0.592 & 71.36 & 3.258 \\
\bottomrule
\end{tabularx}
\par\vspace{4pt}{\footnotesize 每配置 $n=55$；共 495 次检索，均为递归字符分块。相关性以文件名与 PDF 物理页共同匹配。粗体表示九配置中的最优值（延迟取最小），并列最优均标粗。\par}
\endgroup

\tabnote{V=向量，H=混合，H+R=混合加重排。查询时延不含建库；512 档复用已有索引。}
\end{table}
\fig{02_retrieval_quality}{三档分块长度与三种检索方式的质量比较。相同题集与金标页定义下，分别展示 Hit@5、MRR@5、Recall@5；未绘制未经重复实验估计的误差棒。}{1}{fig:retrievalquality}

512+H+R 命中 46/55 题，相对同长度向量检索的 39/55 多命中 7 题，Hit@5 增加 12.73 个百分点；Recall@5 从 59.85\% 增至 72.73\%。重排提高了本轮证据覆盖，但查询时延从 0.1416 秒增至 1.7860 秒。混合本身并不在所有配置和指标上优于向量：512 档两者命中数相同，而混合 Recall 略低，应保留这一结果。

1024+H+R 同样命中 46 题，但 MRR 为 0.5921，Recall 为 71.36\%，平均 3.2577 秒。512 档在命中数相同的情况下排名与覆盖更好，耗时约少 45.2\%。256+H+R 更快，但少命中 4 题。故 512+H+R 是本数据集上合理的后续固定基线；这一判断不是对所有任务的通用最优证明。
\fig{03_retrieval_tradeoff}{检索质量、查询时延与索引规模。时延仅比较查询过程；块数刻画索引规模，不用复用索引的零建库时间制造速度优势。}{1}{fig:tradeoff}
\fig{04_paired_hits}{相同 55 题在九组配置中的逐题命中情况。每格为该题是否至少命中一个金标证据页；分组展示四类任务，便于定位总体均值掩盖的差异。}{0.94}{fig:paired}
\FloatBarrier

\subsection{Agent 效率与工具表现}
\begin{table}[H]\caption{55 题 Agent 基线：执行、工具和效率指标}\label{tab:agent}\centering\small
% Generated by scripts/make_tables.py from data/report_data.json.
\begingroup
\small
\setlength{\tabcolsep}{5pt}
\renewcommand{\arraystretch}{1.16}
\begin{tabularx}{\linewidth}{@{}p{0.29\linewidth}p{0.27\linewidth}>{\raggedright\arraybackslash}X@{}}
\toprule
指标 & 结果 & 分母与判定口径 \\
\midrule
执行成功率 & 53/55（96.36\%） & 后端执行标志；包含失败题。 \\
自动工具集合匹配率 & 46/55（83.64\%） & 与预标允许工具集合精确匹配；不评估参数语义。 \\
AI 辅助语义路由合理率 & 45/55（81.82\%） & AI 同时检查选用工具、调用参数与完整任务意图。 \\
平均端到端延迟 & 64.096 s & $n=55$；包含失败题与保留的原始测量。 \\
P95 端到端延迟 & 168.346 s & $n=55$；按冻结汇总的第 95 百分位口径。 \\
平均执行步骤 & 1.20 步/题 & $n=55$；含直接工具执行与 ReAct。 \\
ReAct 平均执行步骤 & 1.41 步/题 & 仅 $n=27$ 个 ReAct 请求。 \\
完整计量 Token 总量 & 241\,703 token & $n=55$；每题所有 LLM 调用的输入与输出之和。 \\
平均 Token 消耗 & 4\,394.6 token/题 & 分母为 $n=55$ 个计量完整的题目。 \\
Token 计量覆盖率 & 55/55（100.00\%） & 缺失计数不按零处理。 \\
\bottomrule
\end{tabularx}
\par\vspace{4pt}{\footnotesize 自动工具匹配与 AI 语义路由是不同指标。评分未作独立人工复核；执行成功不等于答案正确。\par}
\endgroup

\tabnote{失败题保留在总体分母与计量中；ReAct 步数只对其 27 题求均值。模式间时延差异受题目类型影响，不能当作同题随机对照。}
\end{table}
53/55 的执行成功率为 96.36\%，但执行完成并不能推导内容正确。自动工具集合匹配为 46/55（83.64\%），AI 结合意图及参数审阅后的路由有效为 45/55（81.82\%）；例如 S09 可接受语义上的论文比较，但预标注集合不含该工具，两个口径可能不同。保持原始标注不变，同时保留差异解释。

逐题平均墙钟时延为 64.0958 秒，P95 为 168.3460 秒，平均实际调用 Token 为 4394.6。直接工具路由 28 题平均 102.3914 秒，ReAct 27 题平均 24.3818 秒；对比类任务全部走直接工具路径，题目组成不同，不能据此声称 ReAct 本身必然更快。总体平均步数 1.2，ReAct 子集平均 1.4074，后者才适合描述实际循环推理步数。
\fig{05_agent_performance}{Agent 的逐题效率分布与实际 Token 消耗。保留失败样本，按类别和执行模式展示；每个点表示一个问题，模式对比为观察性描述。}{1}{fig:agent}

\subsection{AI 辅助量表评分}
\subsubsection{评分规则与溯源}
正确性、完整性、引用准确性各按 0/1/2 评分，分别表示核心无效、部分满足、完整满足相应标准。正确性关注关键事实与内部矛盾；完整性关注所问要点覆盖；引用准确性关注实际返回的文件/页能否支持表述。路由采用 0/1，检查工具意图及关键参数。评分参考固定参考答案与真实原文，不只看最终一句是否包含目标数字；执行失败按无有效回答处理。

评分文件记录逐题原因、核对页及不确定性。C01 的“无归纳偏置”绝对化表述被复核为部分正确；C14 虽错误绑定 Vision Mamba，ResNet 部分仍有真实支持，因此三个量表均为 1，而非全部归零。F11 的“最多重试三次/总共尝试三次”在原论文摘要与正文存在表述差异，保留不确定性，不额外据此扣分。原始回答、参考答案和运行元数据没有为获得更好成绩而改写。

\begin{table}[H]\caption{四类问题的 AI 辅助质量评分（各质量维度满分 2）}\label{tab:scores}\centering\small
% Generated by scripts/make_tables.py from data/report_data.json.
\begingroup
\small
\setlength{\tabcolsep}{5pt}
\renewcommand{\arraystretch}{1.16}
\begin{tabularx}{\linewidth}{@{}>{\raggedright\arraybackslash}X r r r r r@{}}
\toprule
问题类别 & $n$ & \shortstack{正确性\\（0--2）} & \shortstack{完整性\\（0--2）} & \shortstack{引用准确性\\（0--2）} & \shortstack{AI 语义路由\\合理率（\%）} \\
\midrule
事实 & 20 & 0.700 & 0.750 & 0.650 & 60.00 \\
对比 & 15 & 1.133 & 0.733 & 1.133 & 100.00 \\
归纳 & 10 & 1.000 & 0.700 & 1.000 & 80.00 \\
推理 & 10 & 1.300 & 1.200 & 1.300 & 100.00 \\
\midrule
\textbf{总体} & \textbf{55} & \textbf{0.982} & \textbf{0.818} & \textbf{0.964} & \textbf{81.82} \\
\bottomrule
\end{tabularx}
\par\vspace{4pt}{\footnotesize 三项答案评分取 0、1、2 分；表中为均值。各维度均覆盖对应类别全部题目，失败题保留。语义路由逐题判 0/1，结果为 AI 辅助评分。\par}
\endgroup

\tabnote{全部 55 题均有评分；“路由”是 AI 语义判断的 0/1 比例。结果尚待独立人工复核，量表均分不是二元答案正确率。}
\end{table}
\fig{06_ai_quality}{AI 辅助评分分布与分任务均值。正确性、完整性、引用准确性均为 0--2 量表；不将其标为独立人工实验或自动正确率。}{1}{fig:ai}

总体正确性、完整性和引用准确性分别为 0.9818/2、0.8182/2、0.9636/2。完整性是三个量表中的最低项。事实类部分题被错误送入元数据工具，导致实验数值或配置缺失；对比类虽工具集合匹配较高，仍会因论文绑定错误或只回答一部分要点而失分。推理类表现相对较好，但样本仅 10 题，不能据此推广到所有推理任务。

\subsection{Bad Case 分析与优化交接}
四类错误采用可重叠标签：检索 10 题、生成 39 题、引用 42 题、Agent 11 题，均来自 AI 辅助诊断。相加为 102 次标签，而非 102 个独立失败题。某题可同时出现范围错配、错误生成与引用不支持；图\ref{fig:badcases}因此采用计数和共现矩阵，不将四类绘成互斥饼图。
\fig{07_bad_cases}{四类 Bad Case 的多标签计数与共现。对角线为各类涉及题数，非对角线为两类共同涉及的题数；标签表达审阅判断，不单独证明模型内部因果。}{1}{fig:badcases}
\begin{table}[H]\caption{代表性问题、可见证据与下一步处理}\label{tab:cases}\centering\small
% Generated by scripts/make_tables.py from data/report_data.json.
\begingroup
\footnotesize
\setlength{\tabcolsep}{5pt}
\renewcommand{\arraystretch}{1.2}
\begin{tabularx}{\linewidth}{@{}p{0.06\linewidth}p{0.15\linewidth}p{0.31\linewidth}>{\raggedright\arraybackslash}X@{}}
\toprule
题号 & 主要问题 & 已观察现象 & 核查证据与待验证改进 \\
\midrule
\textbf{F02} & 生成／传参偏移 & 将集成模型测试集 3.57\% 误写为单模型验证集结果。 & P06，第 2、7 页区分两种设置；建议保留原题限定，并逐项绑定模型、数据划分与指标。 \\
\addlinespace[4pt]
\textbf{F07} & Agent 执行失败 & 事实问题误走元数据工具，字符串拼接 TypeError 导致无答案、无引用。 & P12，第 1 页包含所问设置；建议修正意图判定、规范字段类型，并保留真实失败状态。 \\
\addlinespace[4pt]
\textbf{F16} & 检索证据缺失 & 未返回 ViT 数据集规模所在页，最终拒答三组类别数／图像数。 & P03，第 4 页有完整数据；建议约束目标来源与章节，再检验召回和答案完整性。 \\
\addlinespace[4pt]
\textbf{C14} & 来源匹配／生成 & Vim 匹配成另一篇论文；又将 ResNet 的 ImageNet 指标误归 CIFAR-10。 & P02，第 2、5 页及 ResNet 第 7 页可核对；建议使用确定的题名—文件映射，并按数据集组织证据。 \\
\addlinespace[4pt]
\textbf{S10} & 任务分解／状态 & 三篇论文被交给双篇比较工具；返回错误信息，却以 success=True 结束。 & 三篇原文第 1 页均有相应内容；建议逐篇抽取后汇总，并将工具失败状态向上传递。 \\
\addlinespace[4pt]
\textbf{R08} & 推理归因错误 & 以色彩基线缺少端到端评估，解释复合规则 0/102 不能算独立验证误报率。 & P14，第 8、11 页说明同集设计偏差及冻结后独立验证；建议区分评估对象、设计集与验证集。 \\
\bottomrule
\end{tabularx}
\par\vspace{4pt}{\footnotesize 案例由冻结的 AI 辅助评分与实际结果归纳；P 编号对应文献表，页码均为 PDF 物理页。改进为待实施、待复测建议，不代表已取得优化后结果。\par}
\endgroup

\tabnote{案例均保留原回答及轨迹；原文页出现在 sources 中，只证明页级可见支持，不能还原未记录的全部内部上下文。}
\end{table}

优先改进顺序为：第一，修复 \code{paper_metadata} 的结构化作者字段拼接异常，并将实验结果、训练配置等事实查询路由到证据检索；第二，为论文比较与总结加入显式文献 ID 绑定，防止 Vision Mamba 指向语义分割基准论文；第三，检查回答要点覆盖与数值语境，避免把测试集集成结果改写为验证集单模型结果；第四，引用生成前核验结论与证据的一致性，并对纯诊断性文本设置任务完成检查。

修复应由 C/D 在各自代码中实施，B 使用同一 55 题和固定基线重新运行，保留新版哈希、逐题评分与错误迁移表。当前没有该轮结果，因此本报告没有虚构优化后数值，也没有把本轮长度参数筛选替代课程要求的答案质量优化前后对照。
\FloatBarrier

\section{技术设计文档}
\subsection{评测数据契约与输出结构}
QA 输入以 JSONL 一行一题，核心字段为 \code{id/category/question/reference_answer/gold_evidence/accepted_tool_sets}。gold 包含 \code{source/page/evidence}，文件名与物理页组成严格匹配单位。运行结果补充 \code{success/answer/sources/tools_used/steps/execution_mode/wall_latency_s/token_usage}，原始轨迹保存在嵌套对象中。

每次正式运行创建独立目录，包含参数及指纹元数据、QA 快照、逐题结果、审阅表和汇总。自动指标从逐题结果重新计算；量表评分与其来源另存审计，防止未经评分的旧汇总与新版混用。报告包的 \code{data/report_data.json} 冻结 55 题 Agent、495 条检索、评分与语料记录，图表脚本只消费该快照，不调用模型。

\subsection{可追溯性与续跑设计}
可比性由 QA/语料 SHA-256、后端文件指纹、模型名、包版本与配置共同约束。续跑采用新输出目录，按题 ID 去重复制已落盘结果，仅执行缺失题，并保留 \code{origin_run_id}。索引复用需验证语料、分块与嵌入配置以及已存块的文本/页码/ID，而不是仅检查“目录存在”。对历史片段未保存的模型快照或计量脚本指纹，明确保留不可追溯边界，不事后伪造。

\subsection{运行异常与工程改进记录}
\begin{table}[H]\caption{本地运行异常及已实施的评测工程处理}\small\centering
\begin{tabularx}{\linewidth}{@{}L{27mm}Y Y@{}}\toprule 现象 & 可见证据与处理 & 验证边界\\\midrule
模型缺失 / 显存分配失败 & 预检 Ollama 模型和连接；重用合格索引，避免失败后重复嵌入；固定上下文配置。 & 预热成功只说明可调用，不代表答案正确。\\
Python 原生进程退出 & Windows 事件记录涉及 \code{c10.dll} 访问异常；加无缓冲输出、故障日志、逐题检查点与续跑。 & 不能仅据退出位置认定模型显存是根因。\\
WinError 5 文件替换失败 & 独立临时文件、有限重试与退避，保留耐久元数据快照及修订号；逐题结果刷盘。 & 文件锁故障注入验证写入与恢复，不是模型质量提升。\\\bottomrule\end{tabularx}
\end{table}
这些改进服务于结果完整性与长任务可恢复性，是实际完成的工程稳定性工作；课程要求的生成/Agent 质量改进仍需以修复后的同题复测证明。不得把模拟故障注入的 55 条流程记录当作真实 55 题模型质量数据。

\subsection{图表与报告生成设计}
图采用一致的深蓝、青绿与赭色系统，全部同时导出矢量 PDF 和高分辨率 PNG。使用分面比较、逐题矩阵、分布图、散点图及共现矩阵，分别回答“总体差异、个体差异、效率代价与错误关联”。不使用无重复实验依据的误差棒、显著性标记或推测的优化曲线。表采用三线表、统一单位与小数精度，逐题明细使用跨页续表；图表数值均可从冻结 JSON 重建。

\section{总结与展望}
\subsection{已完成成果与主要认识}
本阶段形成了 15 篇论文、296 页、55 题四类问答的基准，完成真实 Agent 基线、495 条检索记录、AI 辅助评分与四类错误分析，并通过运行来源、哈希及独立目录维持可追溯性。在此样本上，512 字符配合混合检索加重排兼顾页级质量与耗时；Agent 的主要改进空间集中于完整性、文献绑定与引用支持，执行成功和工具集合匹配不足以代表答案可用。

报告交付包括一份可编辑 LaTeX 主文件、八组矢量与位图图形、七个数据表片段、冻结数据及复现脚本。成员 B 的工作价值在于提供能够复查的基线和具体修复入口，而不是以单一高分掩盖失败样本。

\subsection{课程验收前的后续工作}
\begin{table}[H]\caption{待补事项与建议验收产物}\small\centering
\begin{tabularx}{\linewidth}{@{}L{40mm}L{22mm}Y@{}}\toprule 待补工作 & 建议责任 & 可验收产物\\\midrule
独立人工复核 55 题 & B 组织 & 评分者、日期、逐题理由、争议处理与签字记录。\\
路由、异常与论文绑定修复 & C/D 实施，B 复测 & 新版代码指纹、同题新运行、优化前后指标及 Bad Case 迁移。\\
三类分块策略质量比较 & C 实现，B 测量 & 固定语料与题集的 fixed/recursive/semantic 检索质量表。\\
真实界面截图 3--5 张 & A/D & 运行界面及核心区域标注，覆盖实际核心功能。\\
其余模块课程要求 & 对应负责人 & 多格式、嵌入/向量库对照、真实流式等逐项验收材料。\\\bottomrule\end{tabularx}
\end{table}
所有下一轮对照应固定题集与评分规则，失败保留在分母中，旧基线只读保存。若改变语料、模型或金标，应另建实验版本并解释变化，不能直接把不同条件下的得分当作修复收益。
\FloatBarrier

\section{附录}
\subsection{环境与关键配置}
\begin{table}[H]\caption{正式运行记录的环境与配置}\label{tab:env}\centering\small
% Generated by scripts/make_tables.py from data/report_data.json.
\begingroup
\small
\setlength{\tabcolsep}{5pt}
\renewcommand{\arraystretch}{1.16}
\begin{tabularx}{\linewidth}{@{}p{0.16\linewidth}p{0.31\linewidth}>{\raggedright\arraybackslash}X@{}}
\toprule
类别 & 项目 & 版本或配置 \\
\midrule
运行环境 & Python & 3.10.20 \\
 & 操作系统 & Windows（build 26200） \\
\addlinespace[3pt]
生成模型 & Ollama 模型 & \texttt{qwen2.5:7b} \\
\addlinespace[3pt]
检索模型 & Embedding & BAAI/\allowbreak{}bge-\allowbreak{}large-\allowbreak{}zh-\allowbreak{}v1.5 \\
 & Reranker & BAAI/\allowbreak{}bge-\allowbreak{}reranker-\allowbreak{}base \\
 & 推理设备 & \texttt{cpu} \\
\addlinespace[3pt]
依赖版本 & PyTorch & 2.11.0+cu126 \\
 & Sentence Transformers & 6.1.0 \\
 & ChromaDB / rank-bm25 & 1.5.9 / 0.2.2 \\
 & pdfplumber / Ollama SDK & 0.11.10 / 0.6.3 \\
\addlinespace[3pt]
生成参数 & 上下文窗口 & 4\,096 token \\
 & 请求超时 & 120 s \\
\addlinespace[3pt]
Agent 参数 & 分块／重叠 & 512 / 50 字符 \\
\addlinespace[3pt]
检索实验 & 递归字符块长 & 256 / 512 / 1024 字符 \\
\addlinespace[3pt]
检索参数 & 候选／输出 & 20 / 5 个文本块 \\
\addlinespace[3pt]
会话参数 & Memory & 关闭；每题独立会话 \\
\addlinespace[3pt]
评分来源 & 性质 & AI 辅助证据评分；未作独立人工复核 \\
\bottomrule
\end{tabularx}
\par\vspace{4pt}{\footnotesize 模型、依赖与请求参数取自冻结运行元数据；Ollama SDK 版本与独立服务版本不是同一概念。块长单位为字符，上下文窗口单位为 token。\par}
\endgroup

\tabnote{模型名和版本取自运行元数据；CPU 指嵌入与重排设备，不能解释为 Ollama 全链路只用 CPU。}
\end{table}

\subsection{独立报告文件夹与编译}
本目录可以整体复制，编译仅依赖目录内的图表文件与常规 LaTeX 环境；重建图表需要 Python、NumPy 与 Matplotlib，不需要再次调用 Ollama。\code{README.md} 说明脚本入口和字体回退。推荐用 XeLaTeX 编译两次：
\begin{lstlisting}
xelatex -interaction=nonstopmode -halt-on-error main.tex
xelatex -interaction=nonstopmode -halt-on-error main.tex
\end{lstlisting}
报告源文件中的姓名、学号、班级、教师栏位应在提交前填写；填完后重新编译。若使用便携 Tectonic，可执行 \code{tectonic main.tex}。附带 PDF 是本次生成的阅读版本，编辑以 \code{main.tex} 为准。

\subsection{本地实验复现命令}
下列命令在原项目根目录执行，要求使用已安装项目依赖和模型的 NLP 环境。验证与汇总不会重新生成回答；正式运行会调用本地模型并写入新目录。不要覆盖原始基线。
\begin{lstlisting}
# PowerShell: start in the original rag_project directory
$py = "D:\Anaconda\envs\NLP\python.exe"
& $py scripts\module5_eval.py check `
  --qa evaluation\qa.jsonl --data evaluation\corpus

# New Agent run; use a new output directory for each experiment
$stamp = Get-Date -Format "yyyyMMdd_HHmmss"
& $py scripts\module5_eval.py run `
  --mode agent --qa evaluation\qa.jsonl `
  --data evaluation\corpus --chunk-size 512 `
  --out "evaluation\runs\agent55_new_$stamp" --verbose

# Independent retrieval experiment; runs all three retrieval methods
$stamp = Get-Date -Format "yyyyMMdd_HHmmss"
& $py scripts\module5_eval.py run `
  --mode retrieval --qa evaluation\qa.jsonl `
  --data evaluation\corpus --chunk-size 512 `
  --out "evaluation\runs\retrieval55_new_$stamp" --verbose
\end{lstlisting}
256 与 1024 档仅更换 \code{--chunk-size} 并指定新目录。索引复用及续跑的参数和检查由脚本 \code{--help} 给出，不能跳过配置一致性验证。新模型运行不会自动生成独立人工评分；原 AI 评分也不能直接移植给新版回答。

\subsection{原始记录与数据来源}
原项目的主要代码结构如下，按实际文件整理；省略缓存、向量存储与临时日志，不照抄模板中的示例代码树。
\begin{lstlisting}
rag_project/
  app.py
  src/
    data_loader/    # PDF parsing
    chunking/       # chunkers and size statistics
    retrieval/      # Chroma, BM25, hybrid, reranker
    generation/     # RAG pipeline
    agent/          # router, tools, ReAct, memory
    frontend/       # service and UI state
  scripts/
    module5_eval.py
    module5_tokens.py
    module5_charts.py
    run55.py
  evaluation/       # QA, corpus manifest, run records
  docs/
  tests/
  B_Module5_Report_20261014/
\end{lstlisting}
\begin{table}[H]\caption{报告主要数据来源（均在原项目 evaluation 下）}\small\centering
\begin{tabularx}{\linewidth}{@{}L{27mm}Y@{}}\toprule 记录 & 相对目录或文件\\\midrule
QA 与语料 & \code{qa.jsonl}，\code{corpus_manifest.json}\\
Agent 基线 & \code{runs/agent55_20261014_175735_da23d1/}\\
AI 辅助评分 & \code{ai_review_20261014_191222/scores_audit.json}\\
评分后汇总 & \code{ai_review_20261014_191222/summary/summary.json}\\
256 检索 & \code{runs/retrieval55_20261014_210831_0865ca/}\\
512 检索 & \code{runs/retrieval55_20261014_193104_b35edd/}\\
1024 检索 & \code{runs/retrieval55_20261014_212257_280fb9/}\\
配置汇总 & \code{retrieval_comparison_20261014_214329/comparison.json}\\\bottomrule\end{tabularx}
\end{table}
完整 SHA-256 及元数据见本包 \code{data/source_manifest.json}。为保持交付轻量，未复制约百兆字节的原始论文 PDF；参考摘录与页码保留在 QA 中，原论文仍由原项目 \code{evaluation/corpus} 提供。验证原始数据需使用该语料，阅读和编译报告不依赖它。

\subsection{依据文档}
本报告结构与验收边界参考用户提供的《南京农业大学课程实践.docx》与《项目交付模板.docx》；成员比例参考团队分工材料。报告中的所有实验结论以本地冻结记录为依据。八张图均为本次数据绘制或代码依赖示意，不是原论文插图，也不替代模板要求的真实界面截图。

\subsection{文献清单}
以下题名经本地 PDF 第一页文本与评测资料核对，页数为物理页。P 编号同时用于覆盖图；本表不补造未经核验的作者、出版年份或 DOI。
\begingroup\small
% Generated by scripts/make_tables.py from data/report_data.json.
\begingroup
\small
\setlength{\tabcolsep}{5pt}
\renewcommand{\arraystretch}{1.13}
\par\vspace{4pt}{\footnotesize 共 15 篇论文；页数为文件的 PDF 物理页数。文件名与哈希映射保存在包内数据文件。\par}
\begin{longtable}{@{}l >{\raggedright\arraybackslash}p{0.79\linewidth}r@{}}
\caption{15篇本地论文与物理页清单}\label{tab:corpus}\\
\toprule
编号 & 论文题名（原文） & 页数 \\
\midrule
\endfirsthead
\multicolumn{3}{@{}l}{\footnotesize（续表）} \\
\toprule
编号 & 论文题名（原文） & 页数 \\
\midrule
\endhead
\midrule
\multicolumn{3}{r@{}}{\footnotesize 续下页} \\
\endfoot
\bottomrule
\endlastfoot
P01 & Video-\allowbreak{}LLaVA: Learning United Visual Representation by Alignment Before Projection & 14 \\
\addlinespace[3pt]
P02 & Vision Mamba: Efficient Visual Representation Learning with Bidirectional State Space Model & 14 \\
\addlinespace[3pt]
P03 & An Image is Worth 16x16 Words: Transformers for Image Recognition at Scale & 22 \\
\addlinespace[3pt]
P04 & Multi-\allowbreak{}View Semantic–Geometric Perception for Post-\allowbreak{}Disaster Flexible Suction with Metric 3D Scene Reconstruction & 26 \\
\addlinespace[3pt]
P05 & Attention Is All You Need & 15 \\
\addlinespace[3pt]
P06 & Deep Residual Learning for Image Recognition & 9 \\
\addlinespace[3pt]
P07 & SBRF-\allowbreak{}Det: Scale-\allowbreak{}Boundary Representative Fusion for Multi-\allowbreak{}Finding Detection in Panoramic Dental Radiographs & 35 \\
\addlinespace[3pt]
P08 & A Thermal-\allowbreak{}Aware Infrared–Visible Image Fusion Method for Defect Detection in Power Transmission Equipment & 58 \\
\addlinespace[3pt]
P09 & Vision-\allowbreak{}Guided Cone Manipulation with Collaborative Robotic Arm & 10 \\
\addlinespace[3pt]
P10 & Occlusion-\allowbreak{}Aware Watermelon Size Estimation Using YOLO11-\allowbreak{}OBB and Amodal GAN & 8 \\
\addlinespace[3pt]
P11 & How to Choose a Suitable Method in Image Semantic Segmentation—An Experimental Study & 6 \\
\addlinespace[3pt]
P12 & How to Benchmark Vision Foundation Models for Semantic Segmentation? & 10 \\
\addlinespace[3pt]
P13 & An Intelligent Framework for Tooth Segmentation and Classification Through Fusion of Transfer Learning with Advanced Optimization Algorithms Using Biomedical Images & 19 \\
\addlinespace[3pt]
P14 & Zero-\allowbreak{}Shot Water Monitoring from Fixed Drainage Surveillance Imagery Across Day/\allowbreak{}Night and Rain/\allowbreak{}No-\allowbreak{}Rain Conditions: A SAM 3 Benchmark at 25 Urban Sites & 23 \\
\addlinespace[3pt]
P15 & Constructing a Xiangyunsha Defect Dataset and a Lightweight Detection Model for Sustainability & 27 \\
\midrule
\multicolumn{2}{@{}l}{\textbf{合计：15 篇}} & \textbf{296} \\
\end{longtable}
\endgroup

\endgroup

\subsection{55 题评分明细}
三列质量分数依次为正确性、完整性和引用准确性（0--2）；AI 路由为 0/1 的语义判断；自动匹配为固定可接受工具集合匹配。详细问题、参考答案与逐题理由分别见 \code{data/qa.jsonl} 和 \code{data/scores_audit.json}。全部评分仍需独立人工复核。
\begingroup\small
% Generated by scripts/make_tables.py from data/report_data.json.
\begingroup
\small
\setlength{\tabcolsep}{5pt}
\renewcommand{\arraystretch}{1.13}
\par\vspace{4pt}{\footnotesize 答案三维评分为 0--2；AI 语义路由与自动工具匹配为 0/1。AI 语义评分未经独立人工复核。\par}
\begin{longtable}{@{}l l r r r r r@{}}
\caption{55题逐题AI辅助评分与工具匹配明细}\label{tab:questions}\\
\toprule
题号 & 类别 & \shortstack{正确性\\（0--2）} & \shortstack{完整性\\（0--2）} & \shortstack{引用准确性\\（0--2）} & \shortstack{AI 语义路由\\（0/1）} & \shortstack{自动工具匹配\\（0/1）} \\
\midrule
\endfirsthead
\multicolumn{7}{@{}l}{\footnotesize（续表）} \\
\toprule
题号 & 类别 & \shortstack{正确性\\（0--2）} & \shortstack{完整性\\（0--2）} & \shortstack{引用准确性\\（0--2）} & \shortstack{AI 语义路由\\（0/1）} & \shortstack{自动工具匹配\\（0/1）} \\
\midrule
\endhead
\midrule
\multicolumn{7}{r@{}}{\footnotesize 续下页} \\
\endfoot
\bottomrule
\endlastfoot
F01 & 事实 & 2 & 2 & 2 & 1 & 1 \\
F02 & 事实 & 1 & 1 & 1 & 0 & 1 \\
F03 & 事实 & 0 & 0 & 0 & 0 & 0 \\
F04 & 事实 & 0 & 0 & 0 & 0 & 0 \\
F05 & 事实 & 1 & 2 & 1 & 1 & 1 \\
F06 & 事实 & 2 & 2 & 2 & 1 & 1 \\
F07 & 事实 & 0 & 0 & 0 & 0 & 0 \\
F08 & 事实 & 0 & 0 & 0 & 0 & 0 \\
F09 & 事实 & 1 & 1 & 1 & 1 & 1 \\
F10 & 事实 & 2 & 2 & 2 & 1 & 1 \\
F11 & 事实 & 1 & 1 & 1 & 1 & 1 \\
F12 & 事实 & 0 & 0 & 0 & 0 & 0 \\
F13 & 事实 & 0 & 0 & 0 & 0 & 0 \\
F14 & 事实 & 2 & 2 & 2 & 1 & 1 \\
F15 & 事实 & 1 & 1 & 1 & 1 & 1 \\
F16 & 事实 & 0 & 0 & 0 & 1 & 1 \\
F17 & 事实 & 1 & 1 & 0 & 1 & 1 \\
F18 & 事实 & 0 & 0 & 0 & 0 & 1 \\
F19 & 事实 & 0 & 0 & 0 & 1 & 1 \\
F20 & 事实 & 0 & 0 & 0 & 1 & 1 \\
\addlinespace[4pt]
C01 & 对比 & 1 & 1 & 1 & 1 & 1 \\
C02 & 对比 & 1 & 1 & 1 & 1 & 1 \\
C03 & 对比 & 0 & 0 & 0 & 1 & 1 \\
C04 & 对比 & 1 & 1 & 1 & 1 & 1 \\
C05 & 对比 & 1 & 0 & 1 & 1 & 1 \\
C06 & 对比 & 1 & 1 & 1 & 1 & 1 \\
C07 & 对比 & 1 & 1 & 1 & 1 & 1 \\
C08 & 对比 & 1 & 1 & 1 & 1 & 1 \\
C09 & 对比 & 1 & 0 & 1 & 1 & 1 \\
C10 & 对比 & 2 & 0 & 2 & 1 & 1 \\
C11 & 对比 & 1 & 1 & 1 & 1 & 1 \\
C12 & 对比 & 2 & 1 & 2 & 1 & 1 \\
C13 & 对比 & 2 & 1 & 2 & 1 & 1 \\
C14 & 对比 & 1 & 1 & 1 & 1 & 1 \\
C15 & 对比 & 1 & 1 & 1 & 1 & 1 \\
\addlinespace[4pt]
S01 & 归纳 & 1 & 1 & 1 & 1 & 1 \\
S02 & 归纳 & 2 & 1 & 2 & 1 & 1 \\
S03 & 归纳 & 0 & 0 & 0 & 1 & 1 \\
S04 & 归纳 & 1 & 1 & 1 & 1 & 1 \\
S05 & 归纳 & 1 & 0 & 1 & 1 & 1 \\
S06 & 归纳 & 1 & 1 & 1 & 0 & 0 \\
S07 & 归纳 & 1 & 1 & 1 & 1 & 1 \\
S08 & 归纳 & 2 & 1 & 2 & 1 & 1 \\
S09 & 归纳 & 1 & 1 & 1 & 1 & 0 \\
S10 & 归纳 & 0 & 0 & 0 & 0 & 0 \\
\addlinespace[4pt]
R01 & 推理 & 1 & 1 & 0 & 1 & 1 \\
R02 & 推理 & 1 & 1 & 1 & 1 & 1 \\
R03 & 推理 & 2 & 2 & 2 & 1 & 1 \\
R04 & 推理 & 2 & 2 & 2 & 1 & 1 \\
R05 & 推理 & 1 & 1 & 1 & 1 & 1 \\
R06 & 推理 & 1 & 1 & 1 & 1 & 1 \\
R07 & 推理 & 2 & 2 & 2 & 1 & 1 \\
R08 & 推理 & 0 & 0 & 1 & 1 & 1 \\
R09 & 推理 & 1 & 1 & 1 & 1 & 1 \\
R10 & 推理 & 2 & 1 & 2 & 1 & 1 \\
\end{longtable}
\endgroup

\endgroup

\end{document}
